\begin{definition}[A finite witness for the swept-string diagram]
    \label{def:peps_torus_swept_patch_endpoint_geometry}
    \lean{TNLean.PEPS.IsSweptPatchCoordinate,
        TNLean.PEPS.sweptPatchCoordinateDecidable,
        TNLean.PEPS.SweptPatchCoordinate,
        TNLean.PEPS.torusSweptPatchCoordinate,
        TNLean.PEPS.torusSweptPatchWitness}
    \leanok
    Let the horizontal and vertical torus periods be at least seven and six,
    respectively, and fix any translated position $v$. Write
    $p(i,j)=(v_x+i,v_y+j)$, with coordinates taken modulo the periods.
    Complete the displayed four-column, three-row patch
    $\{p(i,j):0\leq i\leq3,\ 0\leq j\leq2\}$ by the side sites
    \[
        p(-1,1),\ p(-1,2),\ p(4,1),\ p(4,2)
    \]
    and the lower sites
    \[
        p(1,-1),\ p(2,-1),\ p(1,-2),\ p(2,-2).
    \]
    Denote this finite set by $R_v$. This is an auxiliary completion of the
    swept-string diagram in \cite[fluxon-braiding passage]{Schuch2010PEPS}.
    The distant lower partner lies beyond the final string detour; it is not
    placed at the endpoint of the displayed lower bond.
\end{definition}

\begin{theorem}[The actual translated twenty-site geometry]
    \label{thm:peps_torus_swept_patch_endpoint_geometry}
    \lean{TNLean.PEPS.card_sweptPatchCoordinate,
        TNLean.PEPS.torusSweptPatchCoordinate_injective,
        TNLean.PEPS.torusSweptPatchWitness_card,
        TNLean.PEPS.torusSweptPatchCoordinate_adj_iff}
    \leanok
    \uses{def:peps_torus_swept_patch_endpoint_geometry}
    The twenty specified vertices are distinct and $|R_v|=20$.
    On their six-column, five-row coordinate envelope, native torus adjacency
    holds exactly when one coordinate is equal and the other differs by one.
    In particular, there are no additional periodic bonds in this envelope.
    These assertions hold at every translated position, including either seam.
\end{theorem}
\begin{proof}\leanok
    All coordinate offsets lie in intervals strictly shorter than their torus
    periods, so equality of their residue classes implies equality of the
    offsets. The same argument applies to unit differences, excluding an
    additional wrap adjacency. Translation preserves both equality and
    adjacency. Count the twelve patch sites, the four side sites, and the four
    lower sites.
\end{proof}

\begin{definition}[Four endpoint loops and partner connectors]
    \label{def:peps_torus_swept_patch_endpoint_walks}
    \lean{TNLean.PEPS.torusSweptPatchEndpoint,
        TNLean.PEPS.torusSweptPatchEndpointLoop,
        TNLean.PEPS.torusSweptPatchConnectorA,
        TNLean.PEPS.torusSweptPatchConnectorB}
    \leanok
    \uses{def:peps_torus_swept_patch_endpoint_geometry}
    Take the four plaquette bases, in order, to be
    \[
        A_2=p(-1,1),\quad A_1=p(3,1),\quad
        B_1=p(1,1),\quad B_2=p(1,-2).
    \]
    At each base take the clockwise four-bond plaquette loop.
    The $A_2$--$A_1$ connector goes up, then right four times, then down.
    The $B_1$--$B_2$ connector goes down three times.
    For the eventual bond assignments, a group element $s$ written along a
    rightward or downward matrix arrow in the source figure corresponds to
    native directed transport $s^{-1}$. The value on an ordered bond is this
    directed transport or its inverse, according to its endpoint order.
\end{definition}

\begin{theorem}[Actual loop and connector containment]
    \label{thm:peps_torus_swept_patch_endpoint_walks}
    \lean{TNLean.PEPS.torusSweptPatchEndpoint_eq,
        TNLean.PEPS.torusSweptPatchEndpointLoop_mem_witness,
        TNLean.PEPS.torusSweptPatchConnectorA_mem_witness,
        TNLean.PEPS.torusSweptPatchConnectorB_mem_witness,
        TNLean.PEPS.torusSweptPatchConnectorA_length,
        TNLean.PEPS.torusSweptPatchConnectorB_length}
    \leanok
    \uses{thm:peps_torus_swept_patch_endpoint_geometry,
        def:peps_torus_swept_patch_endpoint_walks}
    All four endpoint loops and both partner connectors are actual native
    walks contained in $R_v$. The connectors have lengths six and three.
    No adjacency or loop-containment assumption is supplied.
    This establishes only the finite geometry of an explicit four-endpoint
    completion. It does not identify a physical operation with the prescribed
    string braid, or assert parent-Hamiltonian membership.
\end{theorem}
\begin{proof}\leanok
    List the four corners of each endpoint plaquette and use the exact unit
    adjacency relation. The two connector lists consist respectively of seven
    and four vertices of the specified coordinate set. Each successive pair
    differs by one coordinate unit, and every listed vertex belongs to $R_v$.
\end{proof}

\begin{theorem}[Directed transport on the actual partner connectors]
    \label{thm:peps_torus_swept_patch_connector_transport}
    \lean{TNLean.PEPS.regularWalkHolonomy_torusSweptPatchConnectorA,
        TNLean.PEPS.regularWalkHolonomy_torusSweptPatchConnectorB}
    \leanok
    \uses{def:peps_torus_swept_patch_endpoint_walks}
    For any bond assignment, write $R(p)$ and $U(p)$ for its rightward and
    upward directed transports at $p$. The reverse-ordered connector transports
    are
    \[
        H_A=U(p(3,1))^{-1}R(p(2,2))R(p(1,2))R(p(0,2))
        R(p(-1,2))U(p(-1,1))
    \]
    and
    \[
        H_B=U(p(1,-2))^{-1}U(p(1,-1))^{-1}U(p(1,0))^{-1}.
    \]
    These equalities include the ordered-endpoint inversions at either seam.
\end{theorem}
\begin{proof}\leanok
    Expand the six and three successive traversals. Reversing an upward
    traversal replaces its transport by its inverse, and concatenation
    multiplies the successive transports in reverse order.
\end{proof}
